% Generated from locale/tr/content/methods/proofs/resources.tex; do not hand-edit.


\olfileid[tr]{mth}{prf}{res}

\section{Diğer Kaynaklar}

Matematikte kanıtların nasıl yapılacağı üzerine yararlı olabilecek pek
çok kitap vardır. Özellikle \emph{How to Read and do Proofs: An
Introduction to Mathematical Thought Processes} \citep{Solow2013} ile
\emph{How to Prove It: A Structured Approach} \citep{Velleman2019}
kitaplarına bakın. \href{http://www.people.vcu.edu/~rhammack/BookOfProof/BookOfProof.pdf}{\emph{Book
of Proof}} \citep{Hammack2013} ile
\href{https://scholarworks.gvsu.edu/books/7/}{\emph{Mathematical
Reasoning}} \citep{Sandstrum2019}, kanıt üzerine çevrim içi ve ücretsiz
erişilebilen kitaplardır. Felsefeciler
\emph{More Precisely: The Math you need to do Philosophy}
\citep{Steinhart2018} kitabını matematiksel akıl yürütmeye iyi bir giriş
olarak görebilir.

İnternette kanıtlar üzerine çeşitli kısa kılavuzlar da bulunur; örneğin
\href{https://math.berkeley.edu/~hutching/teach/proofs.pdf}{``Introduction
  to Mathematical Arguments''} \citep{Hutchings2003} ve
\href{https://eugeniacheng.com/wp-content/uploads/2017/02/cheng-proofguide.pdf}{``How to
  write proofs''} \citep{Cheng2004}.

\subsection{Güdüleyici Videolar}

Ödevinizi yapmak için hiç isteğiniz yokmuş gibi mi hissediyorsunuz?
Moraliniz mi bozuk? Bu videolar yardımcı olabilir!

\begin{itemize}
\item \url{https://www.youtube.com/watch?v=ZXsQAXx_ao0}
\item \url{https://www.youtube.com/watch?v=BQ4yd2W50No}
\item \url{https://www.youtube.com/watch?v=StTqXEQ2l-Y}
\end{itemize}

